%%
%% Response to Editor — Biodiversity and Conservation (Springer Nature)
%%
\documentclass[12pt]{article}
\usepackage[T1]{fontenc}
\usepackage[utf8]{inputenc}
\usepackage[american]{babel}
\usepackage[margin=2.5cm]{geometry}
\usepackage{setspace}
\onehalfspacing
\usepackage{hyperref}

\pagestyle{empty}

\begin{document}

\noindent
\textbf{Response to Editor}\\[6pt]
Manuscript: \textit{Institutional fragility and climatic exposure drive biocultural vulnerability in traditional agricultural knowledge systems}\\[6pt]
Journal: \textit{Biodiversity and Conservation}\\[6pt]
\today

\vspace{24pt}

\noindent
Dear Editor,

\vspace{12pt}

\noindent
Thank you for considering our manuscript and for your feedback. We have addressed your request as detailed below.

\vspace{18pt}

\noindent
\textbf{Editor's comment:} \textit{Please provide the Email Address of the Corresponding Author in the manuscript.}

\vspace{12pt}

\noindent
\textbf{Response:} The email address of the corresponding author (ldvsantos@uefs.br) has been added to the author field on the title page of the revised manuscript, immediately below the ``[Blinded for peer review]'' designation, in the format:

\vspace{6pt}

\begin{quote}
Corresponding author: \href{mailto:ldvsantos@uefs.br}{ldvsantos@uefs.br}
\end{quote}

\vspace{6pt}

\noindent
This addition ensures that the corresponding author's contact information is available within the manuscript while preserving the blinding of author identities for peer review.

\vspace{36pt}

\noindent
We trust that this revision satisfactorily addresses your request and look forward to the continuation of the review process.

\vspace{36pt}

\noindent
Sincerely,

\vspace{36pt}

\noindent
Luiz Diego Vidal Santos, PhD\\
(Corresponding Author)\\
Universidade Estadual de Feira de Santana (UEFS)\\
E-mail: ldvsantos@uefs.br

\end{document}
